\section{Introduction}
\label{sec:intro}

% P1 — Scene: application-level framing
Sparse convolution is the computational backbone of 3D point-cloud
perception.
LiDAR-based object detection~\cite{centerpoint,second}, semantic
segmentation~\cite{semantickitti}, and scene reconstruction all build on
stacks of sparse convolution layers to process inherently sparse voxelized
representations of 3D sensor data.
In production models such as CenterPoint~\cite{centerpoint} and SECOND~\cite{second}, sparse convolution
accounts for \textbf{TODO}\% of end-to-end inference time, making its
runtime efficiency a first-order concern.
These models operate at extreme sparsity: across a typical detection
network on KITTI, the average per-layer active ratio is below
\textbf{TODO}\%, and even the densest layers
rarely exceed \textbf{TODO}\%.

% P2 — What sparse convolution computes and why it is hard on GPUs
A sparse convolution with kernel size $k$ (typically 3; we use symmetric
$k^3$ kernels throughout) examines $k^3$ neighbor offsets around each
active output voxel.
At each offset, the runtime checks whether the corresponding input
neighbor is also active; if so, the input features at that neighbor are
multiplied against the offset's weight slice and accumulated into the
output.
The set of active neighbors varies per voxel and per offset, producing an
irregular, data-dependent computation pattern.
The central runtime challenge is how to reorganize this irregular work into
the regular, coalesced execution that GPUs require for high utilization.

% P3 — Fundamental tradeoff: skip vs efficiency
Existing systems resolve this challenge by committing to one side of a
fundamental tradeoff between \emph{computation skipping} and
\emph{execution efficiency}.

The \emph{implicit GEMM} approach (SpConv~v2~\cite{spconv2},
TorchSparse++~\cite{torchsparse++}) treats every active output voxel as
if all $k^3$ neighbors were present: it concatenates the full neighborhood
into one large matrix and performs a single fused dense matrix multiply
across all offsets simultaneously.
This yields high GPU utilization and Tensor Core compatibility, but it
cannot skip empty neighbors; missing neighbors simply contribute
zero-valued rows that still consume compute cycles and memory bandwidth.
Profiling a CenterPoint backbone on KITTI reveals that
\textbf{TODO}\% of the fused multiply-accumulate operations
involve such zero-valued inputs, wasted compute that
scales with the kernel volume rather than the actual connectivity.

The \emph{per-offset gather-GEMM-scatter} approach
(MinkowskiEngine~\cite{minkowski}) takes the opposite route: it loops
over the $k^3$ offsets one by one, collects only the active pairs at each
offset, and runs a separate matrix multiply for each.
This avoids wasted computation on empty neighbors.
However, it pays a steep price: up to $k^3$ separate kernel launches per
convolution layer, severe load imbalance across offsets, and either a
naive in-kernel matrix multiply or an extra copy into contiguous buffers
before handing off to a tuned dense GEMM library.
On the same CenterPoint backbone, gather-GEMM-scatter spends
\textbf{TODO}\% of its wall time on launch overhead and data movement
rather than useful arithmetic.

At the low densities that dominate real workloads, neither side works well.
Implicit GEMM wastes overwhelmingly on empty neighbors.
Gather-GEMM-scatter skips the empties, but each per-offset matrix multiply
is too small to saturate the GPU.

% P4 — Our insight: worklist-driven execution
We observe that both sides of this tradeoff organize execution around
\emph{offsets}: implicit GEMM iterates all offsets per output point,
while gather-GEMM-scatter launches one kernel per offset.
Neither organizes execution around the \emph{actual active work}.

\sysname takes a different approach: \emph{worklist-driven execution}.
We precompute all active \texttt{(input, output)} neighbor pairs across
every offset, group them by offset, and flatten the groups into a
contiguous stream of dense GEMM tiles (a worklist).
A single persistent-thread kernel then stride-loops over this worklist.
This design has three key properties:
%
\begin{itemize}
  \item \textbf{Exact work.}
        Offsets with no active pairs produce zero tiles and are skipped
        entirely; computation is strictly proportional to active neighbor
        pairs.
  \item \textbf{Global load balance.}
        All offsets share one kernel launch; the stride loop spreads tiles
        evenly across thread blocks regardless of how unevenly pairs
        distribute across offsets.
  \item \textbf{Dense inner compute.}
        Each tile is a standard dense matrix multiply, fully GPU-friendly
        and Tensor Core compatible.
\end{itemize}

% P5 — Challenge: exact work exposes memory issues
Exact-work execution is not free.
Once convolution is decomposed into per-offset active pairs, the runtime
becomes more fetch-on-demand: input features are gathered indirectly via
pair indices, output accumulation spans multiple offsets (requiring atomic
additions), and tile ordering strongly affects cache reuse.
The challenge is how to recover enough data locality \emph{without}
sacrificing exact work or breaking global load balance.

% P6 — Lightweight locality recovery
\sysname addresses this with three low-complexity dataflow choices, rather
than a heavy scheduling framework:
%
\begin{itemize}
  \item \textbf{Offset-major contiguous tile assignment.}
        Tiles in the worklist are ordered so that consecutive iterations of
        the stride loop tend to share the same kernel offset.
        This improves L2 cache locality for input features while preserving
        the load-balance property of the global stride loop.
  \item \textbf{Weight-fragment persistence.}
        When consecutive tiles share the same offset, the weight slice can
        be kept live in registers instead of being reloaded.
        This turns the offset-major ordering into a direct reduction in
        memory traffic.
  \item We use \textbf{two-level sparsity-aware autotuning} to find the
        optimal tile shape and micro-kernel variant with zero overhead on
        the critical path as runtime sparsity level varies.
        Offline sweeps build both a per-layer config table indexed by
        sparsity bucket and an input-to-layer sparsity mapping; at
        inference time, the runtime reads the input active ratio and
        selects the best config for every layer on the host alone,  
        without any per-layer profiling or device-to-host synchronization.
\end{itemize}
%

% % P7 — Operator-level preview; end-to-end evaluation is in progress
% Table~\ref{tab:intro-perf} previews the operator-level impact of this
% execution model on a controlled microbenchmark.
% On an RTX~2080~Ti with FP32 precision, \sysname outperforms all existing
% sparse convolution systems across the entire active-ratio range tested.
% At 1\% active ratio (typical for LiDAR detection layers), \sysname is
% $5.5\times$ faster than the fastest existing system.
% The advantage holds up to 30\%, where \sysname still maintains a
% $1.8\times$ margin.
% End-to-end model-level evaluation on detection and segmentation benchmarks
% is presented in \S\ref{sec:eval}.

\begin{table}[t]
\centering
\caption{Operator-level kernel latency (ms) for SubM $3{\times}3{\times}3$
convolution, $C_\text{in}{=}C_\text{out}{=}64$, spatial
$128{\times}128{\times}8$, FP32, RTX~2080~Ti.
Speedup is over the fastest existing system at each ratio.
End-to-end model evaluation with real datasets is in
\S\ref{sec:eval} (\emph{in progress}).}
\label{tab:intro-perf}
\small
\begin{tabular}{@{}r rrrr r@{}}
\toprule
Active & \multirow{2}{*}{\sysname} & SpConv & TorchSparse++ & Minkowski- & \multirow{2}{*}{Speedup} \\
ratio  &          & v2     &               & Engine     &  \\
\midrule
0.1\% & \textbf{0.012} & 0.421 & 0.240 & 0.230 & $19.2\times$ \\
0.5\% & \textbf{0.045} & 0.374 & 0.240 & 0.323 & $5.3\times$ \\
1\%   & \textbf{0.045} & 0.403 & 0.246 & 0.379 & $5.5\times$ \\
5\%   & \textbf{0.026} & 0.436 & 0.242 & 0.386 & $9.3\times$ \\
10\%  & \textbf{0.067} & 0.551 & 0.377 & 0.445 & $5.6\times$ \\
20\%  & \textbf{0.240} & 1.251 & 0.685 & 0.908 & $2.9\times$ \\
30\%  & \textbf{0.531} & 1.544 & 0.945 & 1.703 & $1.8\times$ \\
\bottomrule
\end{tabular}
\end{table}

\medskip
\noindent This paper makes the following contributions:
\begin{enumerate}
\item A \textbf{worklist-driven execution model} for sparse convolution
      that achieves computation proportional to active neighbor pairs
      within a single persistent-thread kernel, unifying skip and
      efficiency without the tradeoffs of implicit GEMM or
      gather-GEMM-scatter.

\item \textbf{Lightweight locality recovery} through offset-major tile
      assignment and autotuned weight-stationary micro-kernels, recovering
      data reuse without heavy scheduling.

\item An end-to-end evaluation on 3D object detection and semantic
      segmentation models showing \textbf{TODO}$\times$ speedup over
      the fastest existing sparse convolution runtime, with numerically
      identical outputs.
\end{enumerate}
